\section{Reliability: will it be the same tomorrow, at 3 a.m.?}
\label{sec:reliability}

The clinician needs the same answer tomorrow as today. The sixth question asks
whether the system behaves consistently over time and load, how model drift is
detected and handled, and how uncertainty is quantified and communicated.

[DRAFTING INSTRUCTIONS: in the full-framework, ground reliability in the
credibility assessment of \texttt{asme2018vv40} (ASME V\&V 40) and the uncertainty
quantification in the VVUQ suite carried in the prior review. Describe
reproducibility controls, drift monitoring and revalidation, and confidence
intervals at the point of decision. Reproduce Mermaid
\texttt{adoption/mermaid/diagrams/15-uncertainty-gate-flow.md},
\texttt{16-model-revalidation-gitgraph.md}, and \texttt{17-night-shift-sequence.md}
as colored cfwfig TikZ. Carry the ADOPTION ACTION ``continuous monitoring''. Build
a body-width table of reliability controls and their triggers.]

\begin{cfwfig}
\node[cfone] (a) at (0,0) {Estimate}; \node[cftwo] (b) at (3.2,0) {Confidence interval}; \node[cfdec] (c) at (6.8,0) {Within bound?};
\node[cfhope] (d) at (10.2,0.8) {Report}; \node[cfthree] (e) at (10.2,-0.8) {Escalate};
\draw[cfedge] (a)--(b); \draw[cfedge] (b)--(c); \draw[cfedge] (c)--(d); \draw[cfedge] (c)--(e);
\end{cfwfig}
\vspace{-0.2cm}
\cfwcaption{Figure 7. The uncertainty gate (placeholder; expanded from Mermaid 15
in the full-framework).}

Reliability is answered by honesty about doubt: a quantified confidence interval, a
plan for drift, and a guarantee that 3 a.m. is treated exactly like noon.
